% main_optimization_metropolis.tex
\documentclass[aspectratio=169]{beamer}

% --- Language & encoding ---
\usepackage[utf8]{inputenc}
\usepackage[T1]{fontenc}
\usepackage[english,french]{babel}
\usepackage{lmodern}
\usepackage{microtype}

% --- Math & tables ---
\usepackage{amsmath,amssymb}
\usepackage{booktabs}
\usepackage{graphicx}
\usepackage{changepage}
\usepackage{pgfplots}
\pgfplotsset{compat=1.18} % or your installed version

% --- Metropolis Theme ---
\usetheme[numbering=fraction,progressbar=frametitle]{metropolis}
\setbeamertemplate{frame numbering}[fraction]

% --- Colors (same palette) ---
\definecolor{BootcampBlueDark}{HTML}{1A3C68}
\definecolor{TextBlack}{HTML}{000000}
\definecolor{AccentGold}{HTML}{DAA520}
\definecolor{Crimson}{HTML}{990000}
\definecolor{MatrixGreen}{HTML}{228B22}
\definecolor{MatrixRed}{HTML}{B22222}

\setbeamercolor{block title example}{bg=BootcampBlueDark!90!black,fg=white}
\setbeamercolor{block body example}{bg=BootcampBlueDark!5!white,fg=black}

% === Thick frame title band ===
\setbeamercolor{frametitle}{bg=BootcampBlueDark,fg=white}
\setbeamerfont{frametitle}{series=\bfseries,size=\large}
\setbeamertemplate{frametitle}{
  \nointerlineskip
  \begin{beamercolorbox}[wd=\paperwidth,ht=3.2ex,dp=1.4ex,leftskip=1em,rightskip=1em]{frametitle}
    \usebeamerfont{frametitle}\insertframetitle
  \end{beamercolorbox}
}

% === Custom Definition Block ===
\newenvironment{definitionblock}[1]{%
  \par\vspace{0.4em}%
  \begin{beamercolorbox}[
      wd=\textwidth,
      sep=0.4em,
      leftskip=0.3em,
      rightskip=0.6em,
      rounded=false,
      shadow=false
    ]{block title example}%
    \raisebox{0.15em}{\usebeamerfont{block title example}#1}%
  \end{beamercolorbox}%
  \vspace{-0.4em}%
  \begin{beamercolorbox}[
      wd=\textwidth,
      sep=1em,
      leftskip=0.6em,
      rightskip=0.6em,
      rounded=false,
      shadow=false
    ]{block body example}%
    \usebeamerfont{block body example}%
}{%
  \end{beamercolorbox}%
  \par\vspace{0.6em}%
}

% === Global text formatting ===
\setbeamercolor{normal text}{fg=TextBlack,bg=white}
\setbeamercolor{title}{fg=TextBlack}
\setbeamercolor{structure}{fg=BootcampBlueDark}
\setbeamercolor{progress bar}{fg=BootcampBlueDark,bg=gray!30}

% --- Course Info ---
\title{Lecture 3 : Optimization, Duality and Lagrangian}
\institute{NYU Tandon School, Finance and Risk Engineering}

\metroset{
  sectionpage=progressbar,
  subsectionpage=progressbar
}

% --- Section & subsection transitions ---
\AtBeginSection{
  \frame[plain]{\sectionpage}
}
\AtBeginSubsection{
  \frame[plain]{\subsectionpage}
}

\begin{document}

\maketitle

% ===== Outline =====
\begin{frame}{Outline}
  \tableofcontents[hideallsubsections]
\end{frame}

% ===========================
% SECTION 0 — Introduction
% ===========================
\section{Motivation \& Overview}

\begin{frame}{Outline}
  \tableofcontents[currentsection,hideothersubsections]
\end{frame}

\begin{frame}{Optimization in Finance \& Economics}

\begin{itemize}
  \item \textbf{Central paradigm:} nearly every decision problem can be written as
  \[
  \min_{x\in\mathcal{X}} f(x)
  \quad \text{subject to constraints (if any).}
  \]
  where \(f(x)\) represents a cost, loss, or risk,  
  and \(\mathcal{X}\) is the feasible set (budgets, regulations, etc.).
\pause
  \vspace{0.5em}

  \item \textbf{Portfolio theory (Markowitz):}  
  choose weights \(w\) to minimize variance for a given expected return:
  \[
  \min_{w}\; w^\top \Sigma w
  \quad \text{s.t. } \mathbf{1}^\top w = 1,\ \mu^\top w = \bar{r}.
  \]

  \vspace{0.5em}
\pause
  \item \textbf{Utility maximization:}  
  maximize expected utility \(E[U(W)]\) under risk constraints.

  \vspace{0.5em}

  \item \textbf{Pricing \& hedging:}  
  option replication and risk-minimizing portfolios often come from minimizing expected squared hedging errors.

\end{itemize}

\end{frame}

% ------------------------------------------------------

\begin{frame}{Optimization Across Quantitative Tasks(1)}

\textbf{Estimation \& Calibration as Optimization:}

\begin{itemize}
  \item \textbf{Least Squares / MLE:}
  \[
  \min_{\theta} \sum_{i}(y_i - f(x_i;\theta))^2
  \quad \text{or} \quad
  \max_{\theta} \ell(\theta)=\log p(y|\theta).
  \]
  Estimation problems become unconstrained or constrained optimizations.

  \vspace{0.5em}
\pause
  \item \textbf{Risk management:}  
  Minimize a risk functional, let \(L(x)\) be the portfolio loss.
  \[
  \min_{x} \text{CVaR}_\alpha(L(x))
  \quad \text{or} \quad
  \min_x \text{Var}(L(x))  
  .\]

\end{itemize}

\end{frame}

\begin{frame}{Optimization Across Quantitative Tasks (2)}

\textbf{Estimation \& Calibration as Optimization:}

\begin{itemize}

  \item \textbf{Dynamic optimization:}  
  Bellman/HJB equations in continuous time are optimization principles:
  \[
  V(t,x)=\sup_{u(\cdot)} E\!\left[\int_t^T U(X_s,u_s)\,ds + G(X_T)\right].
  \]
  Optimal control → optimal policies.
\pause
  \vspace{0.5em}

  \item \textbf{Machine learning \& calibration:}  
  fitting volatility surfaces, neural networks, or pricing models  
  all rely on gradient-based optimization of loss functions.

\end{itemize}

\end{frame}
% ------------------------------------------------------

\begin{frame}{Economic Interpretation of Optimization}

\begin{itemize}
  \item \textbf{Convexity:}  
        ensures stability and uniqueness of equilibria and prices.
\pause
  \item \textbf{Sensitivity:}  
        derivatives of optimal value w.r.t. parameters give comparative statics and risk sensitivities (the “Greeks’’ in finance).
\pause
  \item \textbf{Broader scope:}  
        optimization unifies:
        \begin{itemize}
          \item decision theory (utility maximization),
          \item econometrics (likelihood optimization),
          \item risk control (variance/CVaR minimization),
          \item and dynamic programming.
        \end{itemize}
\end{itemize}

\end{frame}

% ======================================
% SECTION 1 — Unconstrained Optimization
% ======================================
\section{Unconstrained Optimization}

\begin{frame}{Outline}
  \tableofcontents[currentsection, hideothersubsections]
\end{frame}

% ======================================================
\subsection{First-Order Conditions}
% ======================================================

\begin{frame}{Stationary Points \& First-Order Conditions}

\begin{definitionblock}{First-Order Necessary Condition (FONC)}
Let \(f:\mathbb{R}^n\to\mathbb{R}\) be differentiable.  
If \(x^\star\) is a local minimum (or maximum), then
\vspace{-4mm}
\[
\nabla f(x^\star) = 0.
\vspace{-5mm}
\]
\vspace{-4mm}
\end{definitionblock}
\vspace{-4mm}
\pause
\begin{itemize}
  \item The gradient gives the direction of steepest ascent:
        \(\nabla f(x) = 0\)  means there is no ascent direction left.
        \pause
  \item Any \textbf{stationary point} (where gradient vanishes) is a \textit{candidate}
        for an extremum.
        \pause
  \item To classify the point (min, max, saddle), we need \textbf{second-order information}
        from the Hessian matrix.
        \pause
\end{itemize}
\pause
\vspace{0.4em}
\textbf{Geometric idea:}  
At an optimum, the tangent plane is horizontal (slope = 0 in all directions).

\end{frame}

\begin{frame}{Example 1: One-Variable Case}

\textbf{Function: } \(f(x) = x^3 - 3x^2 + 2.\)
\pause
\[
f'(x) = 3x^2 - 6x = 3x(x-2), \quad f''(x) = 6x - 6.
\]

\textbf{Stationary points: } \(x_1=0,\ x_2=2.\)
\pause
\begin{itemize}
  \item At \(x_1=0\): \(f''(0) = -6 < 0\) → local \textbf{maximum}.
  \item At \(x_2=2\): \(f''(2) = 6 > 0\) → local \textbf{minimum}.
\end{itemize}
\pause
\textbf{Interpretation:}  
Slope is zero at both points, but curvature (sign of \(f''\)) distinguishes them.

\vspace{0.5em}
\textcolor{BootcampBlueDark}{Graphically:} derivative crosses 0 → switch from ↑ to ↓ or ↓ to ↑.

\[
\text{Local min if } f' \text{ changes from negative to positive.}
\]

\end{frame}

\begin{frame}{Example 2: Multivariable Case}

\textbf{Function: } \(f(x,y) = x^2 + xy + y^2 - 4x - 2y.\)
\vspace{-2mm}
\pause
\[
\nabla f(x,y) =
\begin{pmatrix}
2x + y - 4 \\[2pt]
x + 2y - 2
\end{pmatrix}.
\vspace{-3mm}
\]
\vspace{-3mm}
\textbf{Stationary point:}
\[
\nabla f(x^\star,y^\star)=0
\Rightarrow
\begin{cases}
2x^\star + y^\star = 4,\\
x^\star + 2y^\star = 2.
\end{cases}
\Rightarrow
(x^\star,y^\star) = \left(\tfrac{2}{3}, \tfrac{2}{3}\right).
\]
\pause
Compute Hessian:
$
\nabla^2 f(x,y)=
\begin{pmatrix}
2 & 1\\[2pt]
1 & 2
\end{pmatrix}.
$

\begin{itemize}
  \item Eigenvalues \(= 1,3 > 0\) so the Hessian is positive definite.
  \item So \((\tfrac{2}{3},\tfrac{2}{3})\) is a \textbf{local (and global) minimum.}
\end{itemize}
\pause
\vspace{0.4em}
\textbf{Geometric interpretation:}  
The level sets of \(f\) are ellipses; the minimum lies at the center of symmetry.
\pause
\end{frame}

\begin{frame}{Counterexample: Stationary Point $\neq$ Extremum (1)}

\textbf{Goal:} Show that $\nabla f(x^\star)=0$ does not always imply a minimum or maximum.
\pause
\begin{definitionblock}{Example (Saddle Point)}
Consider \(f(x,y) = x^2 - y^2.\)
$
\nabla f(x,y) =
\begin{pmatrix}
2x \\[2pt]
-2y
\end{pmatrix},
$ and $
\nabla^2 f(x,y) =
\begin{pmatrix}
2 & 0 \\[2pt]
0 & -2
\end{pmatrix}.
$
\vspace{-2mm}
\end{definitionblock}
\pause
\vspace{-4mm}
\begin{itemize}
  \item \(\nabla f(0,0)=0\) $\Rightarrow$ \((0,0)\) is a stationary point.
  \pause
  \item But the Hessian is \textbf{indefinite} (one $+$ and one $-$ eigenvalue).  \pause
        \(\Rightarrow\) neither min nor max: it’s a \textbf{saddle point.}
        \pause
  \item Along \(x\)-axis: \(f(x,0)=x^2\) → convex → looks like a minimum.  
  \item Along \(y\)-axis: \(f(0,y)=-y^2\) → concave → looks like a maximum.
\end{itemize}

\textbf{Interpretation:} function curves upward in one direction, downward in another.

\textbf{Graph:} saddle-shaped surface — like a horse saddle or mountain pass.

\end{frame}

\begin{frame}{\textbf{Graph:} saddle-shaped surface — like a horse saddle or mountain pass.}

\begin{center}
\begin{tikzpicture}
\begin{axis}[
  view={55}{35},
  domain=-2:2, y domain=-2:2,
  samples=35,
  axis lines=middle,
  xlabel={$x$}, ylabel={$y$}, zlabel={$f(x,y)$},
  title={Saddle surface $z=x^2-y^2$},
  colormap/cool,
  colorbar,
  grid=both,
  zmax=4, zmin=-4,
]
  \addplot3[surf, shader=interp, opacity=0.9] {x^2 - y^2};
\end{axis}
\end{tikzpicture}
\vspace{-2mm}
\end{center}
\vspace{-5mm}
\small The surface curves up in the $x$-direction and down in the $y$-direction, it is what we call a saddle point.
\end{frame}

% ======================================================
\subsection{Second-Order Conditions}
% ======================================================

\begin{frame}{Curvature \& Classification}

\begin{definitionblock}{Second-Order Test (Unconstrained, $C^2$)}
Suppose \(f:\mathbb{R}^n\to\mathbb{R}\) is \(C^2\) and \(\nabla f(x^\star)=0\).
\begin{itemize}
  \item If \(\nabla^2 f(x^\star) \succ 0\) (positive definite) \(\Rightarrow\) \textbf{strict local minimum}.
  \item If \(\nabla^2 f(x^\star) \prec 0\) (negative definite) \(\Rightarrow\) \textbf{strict local maximum}.
  \item If \(\nabla^2 f(x^\star)\) is \textbf{indefinite} \(\Rightarrow\) \textbf{saddle point}.
\end{itemize}
\end{definitionblock}
\pause
\begin{itemize}
  \item Interpretation via Taylor 2nd-order:
  \[
  f(x^\star+h) \approx f(x^\star) + \tfrac12\,h^\top \nabla^2 f(x^\star)\,h.
  \]
  \item Sign of the quadratic form decides local geometry.
\end{itemize}

\end{frame}

\begin{frame}{How to Check Definiteness in Practice}
\textbf{Eigenvalue test.}
\begin{itemize}
  \item All eigenvalues \(>0\) \(\Rightarrow\) PD.
  \pause
  \item All eigenvalues \(<0\) \(\Rightarrow\) ND.
  \pause
  \item Mixed signs \(\Rightarrow\) Indefinite.
  \pause
\end{itemize}
\pause

\medskip
\textbf{Semidefinite (PSD/NSD)}: inconclusive for strict classification.
\begin{itemize}
  \item \(\nabla^2 f(x^\star)\succeq 0\) (PSD) with some zero eigenvalues \(\Rightarrow\) need higher-order terms.
\end{itemize}

\end{frame}

\begin{frame}{Example (1D): Classify with $f''$}

\textbf{Function: } \(f(x)=x^4.\)

\[
f'(x)=4x^3,\quad f''(x)=12x^2.
\]
\pause
\begin{itemize}
  \item Stationary point: \(x^\star=0\) since \(f'(0)=0\).
  \pause
  \item \(f''(0)=0\) (semidefinite case) \(\Rightarrow\) second-order test \textbf{inconclusive}.
  \pause
\end{itemize}
\pause
Check higher order:
\[
f(x)=x^4 \ge 0,\quad \text{and } f(x)>0 \text{ for } x\neq 0,
\]
so \(x^\star=0\) is actually a \textbf{(non-strict) local minimum} but not detected by \(f''\) alone.
\pause
\textcolor{BootcampBlueDark}{Takeaway:} If \(f''(x^\star)=0\), examine higher derivatives or compare values directly.

\end{frame}

\begin{frame}{Example (Multi-D): Quadratic Form \& Classification}

Consider
\[
f(x)=\tfrac12\,x^\top Q x + q^\top x + c, \qquad x\in\mathbb{R}^n.
\vspace{-8mm}
\]

\[
\nabla f(x)=Qx+q,\qquad \nabla^2 f(x)=Q \quad(\text{constant}).
\]
\pause
\textbf{Stationary point (if invertible):}
\[
x^\star = -Q^{-1}q.
\]
\pause
\begin{itemize}
  \item If \(Q\succ 0\) (PD) \(\Rightarrow\) \textbf{strict global minimum} at \(x^\star\).
  \item If \(Q\prec 0\) (ND) \(\Rightarrow\) \textbf{strict global maximum} at \(x^\star\).
  \item If \(Q\) indefinite \(\Rightarrow\) \textbf{saddle}.
\end{itemize}

\textbf{Example:}
\[
Q=\begin{pmatrix}2&1\\1&2\end{pmatrix}\ (\lambda=1,3>0)\Rightarrow \text{PD} \Rightarrow \text{strict min}.
\]

\end{frame}

\begin{frame}{Example (Multi-D): Inconclusive Semidefinite Case}

\textbf{Function: } \(f(x,y)=x^4+y^2.\)

\[
\nabla f(x,y)=\begin{pmatrix}4x^3\\2y\end{pmatrix},
\quad
\nabla^2 f(x,y)=
\begin{pmatrix}
12x^2 & 0\\
0 & 2
\end{pmatrix}.
\]
\pause
At \((0,0)\):
\[
\nabla f(0,0)=0,\qquad
\nabla^2 f(0,0)=\begin{pmatrix}0&0\\0&2\end{pmatrix}\ \text{(PSD, not PD)}.
\]
\pause
\begin{itemize}
  \item SOC inconclusive (one zero eigenvalue).
  \item But \(f(x,y)=x^4+y^2\ge 0\) and \(=0\) only at \((0,0)\) \(\Rightarrow\) \textbf{strict local minimum}.
\end{itemize}

\textcolor{BootcampBlueDark}{Moral:} if Hessian is only semidefinite, check higher-order terms or compare values.

\end{frame}

\begin{frame}{Decision Tree (Unconstrained, $C^2$)}

\begin{itemize}
  \item Step 1: Find stationary points: solve \(\nabla f(x)=0\).
  \pause
  \item Step 2: Evaluate \(H=\nabla^2 f(x^\star)\).
  \pause
  \item Step 3: Classify:
  \begin{itemize}
     \item \(H\succ 0\) \(\Rightarrow\) strict local minimum.
     \item \(H\prec 0\) \(\Rightarrow\) strict local maximum.
     \item \(H\) indefinite \(\Rightarrow\) saddle.
     \item \(H\succeq 0\) or \(H\preceq 0\) but singular \(\Rightarrow\) \textbf{inconclusive}: inspect higher-order terms or compare function values.
     \pause
  \end{itemize}
\end{itemize}

\end{frame}
% ======================================
% SECTION 2 — Convexity Tools
% ======================================
% ======================================================
\section{Convexity \& Concavity}
% ======================================================

\begin{frame}{Outline}
  \tableofcontents[currentsection, hideothersubsections]
\end{frame}

\begin{frame}{Convex Functions: Definition \& Geometry}

\begin{definitionblock}{Convexity}
A function \(f:\mathbb{R}^n\to\mathbb{R}\) is \textbf{convex} on a convex domain \(D\) if
\[
f(\lambda x+(1-\lambda)y)
\le
\lambda f(x)+(1-\lambda) f(y)
\quad\forall x,y\in D,\;\lambda\in[0,1].
\vspace{-5mm}
\]
\end{definitionblock}
\pause
\begin{itemize}
  \item Graphically: the line segment joining \((x,f(x))\) and \((y,f(y))\) lies \emph{above} the graph of \(f\).
  \item The \textbf{epigraph} \(\mathrm{epi}(f)=\{(x,t):t\ge f(x)\}\) is convex.
  \item Convex functions have no “local traps”: every local minimum is global.
\end{itemize}


\end{frame}

\begin{frame}{Concavity \& Linear Examples}

\begin{definitionblock}{Concavity}
A function \(f\) is \textbf{concave} on \(D\) if
\[
f(\lambda x+(1-\lambda)y)
\ge
\lambda f(x)+(1-\lambda)f(y)
\quad\forall\lambda\in[0,1].
\]
\vspace{-5mm}
\end{definitionblock}
\pause
\vspace{-5mm}
\begin{itemize}
  \item Equivalent to convexity of \(-f\).
  \item Common examples:
    \begin{itemize}
      \item \(f(x)=\log x\), \(x>0\): concave.
      \item \(f(x)=e^x\): convex.
      \item \(f(x)=x^2\): convex.
      \item \(f(x)=-x^2\): concave.
    \end{itemize}
  \item In economics and finance:
    \begin{itemize}
      \item Utility functions are typically concave (risk aversion).
      \item Cost or loss functions are convex (penalize deviations).
    \end{itemize}
\end{itemize}
\end{frame}

\begin{frame}{First-Order Characterization (Gradient Form)}

If \(f\) is differentiable, then \(f\) is convex on a convex set \(D\)
\textbf{iff}:
\[
f(y) \ge f(x) + \nabla f(x)^\top (y-x)
\quad\forall x,y\in D.
\]
\pause
\begin{itemize}
  \item The tangent hyperplane at \(x\) lies \textbf{below} the graph of \(f\).
  \item The gradient \(\nabla f(x)\) provides a global supporting hyperplane.
  \item Equality for all \(x,y\) $\iff$ \(f\) is affine.
\end{itemize}


\end{frame}

\begin{frame}{Order on Symmetric Matrices}

\begin{definitionblock}{Loewner Order}
Let \(A,B\in\mathbb{S}^n\) (real symmetric matrices).
We say that
\vspace{-3mm}
\[
A \preceq B
\quad\Longleftrightarrow\quad
B-A \text{ is positive semidefinite (PSD)}.
\vspace{-3mm}
\]
Similarly,
\vspace{-3mm}
\[
A \prec B
\quad\Longleftrightarrow\quad
B-A \text{ is positive definite (PD)}.
\]
\end{definitionblock}
\pause
\begin{itemize}
  \item \(A\succeq 0\): \(A\) is \textbf{positive semidefinite} (PSD).
  \item \(A\succ 0\): \(A\) is \textbf{positive definite} (PD).
  \item Quadratic form characterization:
  \[
  A\succeq 0 \;\Longleftrightarrow\; x^\top A x \ge 0,\ \forall x\in\mathbb{R}^n.
  \]
  \item Spectral characterization:
  \[
  A\succeq 0 \;\Longleftrightarrow\; \lambda_i(A)\ge 0\ \text{ for all eigenvalues } \lambda_i.
  \]
\end{itemize}

\textcolor{BootcampBlueDark}{Intuition:}
the order \( \preceq \) compares the \textbf{curvatures} represented by matrices.

\end{frame}

\begin{frame}{Second-Order Convexity Test (Hessian Criterion)}

\begin{definitionblock}{Second-Order Test ($C^2$ case)}
If \(f\in C^2(D)\) and \(\nabla^2 f(x)\succeq 0\) for all \(x\in D\),
then \(f\) is \textbf{convex}.
If \(\nabla^2 f(x)\succ 0\) for all \(x\), \(f\) is \textbf{strictly convex}.
\end{definitionblock}
\pause
\begin{itemize}
  \item Positive semidefinite (PSD) Hessian: curvature never bends downward.
  \item Strict PD Hessian: curvature strictly positive everywhere $\Rightarrow$ unique minimizer.
  \item In 1D: \(f''(x)\ge 0 \Rightarrow f\) convex.
\end{itemize}
\pause
\textcolor{BootcampBlueDark}{Example:}
\[
f(x_1,x_2)=x_1^2+x_2^2, \quad
\nabla^2 f(x)=\begin{pmatrix}2&0\\0&2\end{pmatrix}\succ0.
\Rightarrow f\text{ strictly convex}.
\]

\end{frame}

\begin{frame}{Strong Convexity \& Lipschitz Gradients}

\textbf{Strong convexity:}
\(f\) is $\mu$-\textbf{strongly convex} if
\[
f(y)\ge f(x)+\nabla f(x)^\top(y-x)+\frac{\mu}{2}\|y-x\|^2,
\quad \mu>0.
\]
\pause
\begin{itemize}
  \item $\Rightarrow$ Hessian uniformly bounded: \(\nabla^2 f(x)\succeq \mu I\).
  \item Ensures uniqueness of minimizer and faster convergence of gradient descent.
\end{itemize}
\pause
\textbf{Lipschitz gradient:}
\(f\) has $L$-Lipschitz gradient if
\[
\|\nabla f(x)-\nabla f(y)\|\le L\|x-y\|.
\]
\pause
\begin{itemize}
  \item $\Rightarrow$ Hessian $\preceq L I$ (bounded curvature).
  \item Together, $\mu I\preceq \nabla^2 f(x)\preceq L I$ defines \textbf{well-conditioned} problems.
\end{itemize}

\end{frame}

\begin{frame}{Convexity Properties \& Operations}

\textbf{Preservation rules:}
\begin{itemize}
  \item Nonnegative weighted sums of convex functions are convex:
  \[
  f(x)=\sum_i \alpha_i f_i(x),\ \alpha_i\ge0.
  \]
  \item Composition with affine map: \(f(Ax+b)\) convex if \(f\) convex.
  \item Pointwise supremum of convex functions is convex.
\end{itemize}
\pause
\textbf{Common convex functions:}
\[
\begin{aligned}
&\|x\|_2,\ \|x\|_1,\ \|x\|_\infty,\\
&x\mapsto e^x,\ \log\!\sum_i e^{x_i}\ (\text{log-sum-exp}),\\
\vspace{-3mm}
&x\mapsto (x_+)^2,\quad x_+\!:=\max(0,x).
\end{aligned}
\]
\pause
\end{frame}

\begin{frame}{Why Convexity Matters in Optimization}

\begin{block}{Fundamental Principle}
\textbf{For a convex function on a convex feasible set:}  
\[
\text{Every local minimum is a global minimum.}
\]
\end{block}

\begin{itemize}
  \item \textbf{Convex objective:} \(f(\lambda x + (1-\lambda)y) \le \lambda f(x) + (1-\lambda)f(y)\).
  \item \textbf{Convex constraints:} feasible region is a convex set.
  \item \textbf{Implication:} optimization landscape has \emph{no spurious minima}, no traps, no bad basins.
\end{itemize}

\vspace{0.5em}

\textcolor{BootcampBlueDark}{\textbf{Consequences for algorithms:}}
\begin{itemize}
  \item Gradient descent is guaranteed to converge (under mild assumptions).
  \item No need for multistart strategies or global search routines.
  \item Uniqueness of the solution when \(f\) is strictly convex.
\end{itemize}


\end{frame}

\begin{frame}{Concavity in Practice (Economic Interpretation)}

\textbf{Concave functions:} $f$ is concave $\iff$ $-f$ convex.

\begin{itemize}
  \item \textbf{Utility:} risk-averse agent $\Rightarrow$ concave utility $U(w)$, decreasing marginal utility.
  \item \textbf{Production:} diminishing returns $\Rightarrow$ concave production function.
  \item \textbf{Log-likelihood:} many statistical log-likelihoods are concave $\Rightarrow$ easy maximization.
\end{itemize}

\textcolor{BootcampBlueDark}{Example:}
\[
U(w)=\sqrt{w}\quad\Rightarrow\quad
U''(w)=-\frac{1}{4w^{3/2}}<0 \text{ (concave)}.
\]

\textbf{Implication:}
Maximization of concave functions behaves like minimization of convex functions:
\[
\nabla^2 f(x)\preceq 0 \;\Rightarrow\; \text{concave}.
\]

\end{frame}

% ======================================
% SECTION 3 — Constrained Optimization
% ======================================
\section{Constrained Optimization \& Lagrange Multipliers}

\begin{frame}{Outline}
  \tableofcontents[currentsection, hideothersubsections]
\end{frame}

\subsection{Equality Constraints}
\begin{frame}{Lagrange Multipliers: Motivation \& Idea}

\textbf{Goal:} minimize \(f(x)\) subject to a constraint \(g(x)=c\).
\pause
\begin{itemize}
  \item The feasible set is a \textbf{curve} or \textbf{surface} defined by \(g(x)=c\).
  \item At the optimum, you cannot move along the constraint to decrease \(f\) further.
  \item Geometric insight:
        \[
        \nabla f(x^\star)\ \text{is parallel to}\ \nabla g(x^\star).
        \]
  \item So for some scalar \(\lambda^\star\):
        \[
        \nabla f(x^\star) = \lambda^\star \nabla g(x^\star).
        \]
\end{itemize}

\end{frame}

\begin{frame}{Lagrange Multipliers: Motivation \& Idea}
\begin{center}
\begin{tikzpicture}[scale=1.1]

  % Axes
  \draw[->] (-0.2,0) -- (4.5,0) node[right] {$x_1$};
  \draw[->] (0,-0.2) -- (0,3.3) node[above] {$x_2$};

  % Point optimum
  \coordinate (xstar) at (2,1.6);

  % Contrainte g(x) = c (bleu) : passe par xstar
  \draw[thick, BootcampBlueDark]
    (0.5,0.6)
      .. controls (1.0,2.7) and (1.5,1.6) ..
    (xstar)
      .. controls (2.5,1.6) and (3.2,2.6) ..
    (4.0,2.2);
  \node[BootcampBlueDark] at (3.7,2.8) {$g(x)=c$};

  % Niveau de f(x) = const (or) : tangent à g en xstar
  \draw[thick, AccentGold]
    (0.6,1.3)
      .. controls (1.2,0.8) and (1.5,1.6) ..
    (xstar)
      .. controls (2.5,1.6) and (3.1,0.9) ..
    (3.8,1.0);
  \node[AccentGold] at (3.7,0.7) {$f(x)=\text{const}$};

  % Point optimum
  \fill[AccentGold] (xstar) circle (2pt);
  \node[above right] at (xstar) {$x^\star$};

  % Gradients en xstar (colinéaires et normaux aux courbes)
  \draw[->, thick, BootcampBlueDark]
    (xstar) -- ++(0,1.0) node[above] {$\nabla g(x^\star)$};
  \draw[->, thick, AccentGold]
    (xstar) -- ++(0,-1.0) node[below] {$\nabla f(x^\star)$};

\end{tikzpicture}
\end{center}

\small Level sets of \(f\) touch the constraint surface at the optimum.
\end{frame}

\begin{frame}{Lagrange Multipliers: Formalism}

\begin{definitionblock}{Lagrangian Function}
Consider the constrained problem
\vspace{-5mm}
\[
\min_x f(x) \quad \text{s.t. } g(x)=c.
\vspace{-5mm}
\]
\pause
We define the Lagrangian as
\vspace{-5mm}
\[
\mathcal{L}(x,\lambda)=f(x) + \lambda \left(g(x)-c\right).
\vspace{-5mm}
\]
\vspace{-5mm}
\end{definitionblock}
\vspace{-5mm}
\pause
\textbf{First-order conditions (FOC):}
\[
\nabla_x \mathcal{L}(x^\star,\lambda^\star)=
\nabla f(x^\star)+\lambda^\star\nabla g(x^\star)=0,
 \text{ and }
g(x^\star)=c.
\]

\begin{itemize}
  \item The multiplier \(\lambda^\star\) is the \textbf{shadow price}:
    \[
    \lambda^\star = \frac{\partial f^\star}{\partial c}.
    \]
  \item Economic meaning: marginal cost/value of relaxing the constraint.
  \item In finance: value of relaxing budget, leverage, variance, or risk constraints.
\end{itemize}

\end{frame}

\begin{frame}{Example 1 (1D): Minimizing a Quadratic}

\textbf{Problem:}
\[
\min_{x} \; f(x)=x^2 
\quad\text{s.t.}\quad g(x)=x-1=0.
\]
\pause
\textbf{Lagrangian:}
\[
\mathcal{L}(x,\lambda)=x^2+\lambda(x-1).
\]
\pause
\textbf{FOC:}
\[
\partial_x \mathcal{L}=2x+\lambda=0,
\qquad x-1=0.
\]
\pause
Solve:
\[
x^\star=1,\qquad \lambda^\star=-2.
\]

\textbf{Interpretation:} Increasing the constraint value \(c=1\) slightly would increase the optimal value \(f^\star=c^2\);  
\(\lambda^\star = f'(1) = -2.\)

\end{frame}

\begin{frame}{Example 2 (Multivariate): Minimize a Quadratic on a Line}

\textbf{Problem:}
\[
\min_{x\in\mathbb{R}^2} \; f(x)=x_1^2+x_2^2
\quad\text{s.t.}\quad g(x)=x_1+x_2-1=0.
\]
\pause
\textbf{Lagrangian:}
\[
\mathcal{L}(x,\lambda)=x_1^2+x_2^2+\lambda(x_1+x_2-1).
\]
\pause
\textbf{FOC:}
\[
\begin{cases}
2x_1+\lambda=0,\\
2x_2+\lambda=0,\\
x_1+x_2=1.
\end{cases}
\]
\pause
From first two: $
x_1=x_2.
$
Then constraint gives:
\[
x_1=x_2=\tfrac{1}{2},\qquad
\lambda^\star=-1.
\]

\textbf{Conclusion:} The closest point on the line \(x_1+x_2=1\) to the origin is \((\tfrac{1}{2},\tfrac{1}{2})\).

\end{frame}

% \begin{frame}{Interpretation of the Lagrange Multiplier}

% The Lagrange multiplier is the \textbf{shadow price} of the constraint.

% \begin{block}{Shadow Price Interpretation}
% \[
% \lambda^\star = \frac{\partial f^\star}{\partial c}
% \quad\text{(marginal cost/value of relaxing the constraint).}
% \]
% \end{block}

% \begin{itemize}
%   \item \textbf{Budget constraint:} in utility maximization, \(\lambda^\star\) is the marginal utility of wealth.
%   \item \textbf{Risk constraint:} in portfolios, \(\lambda^\star\) is the value of relaxing a variance or VaR bound.
%   \item \textbf{Calibration:} \(\lambda^\star\) measures sensitivity of the optimal fit to constraint adjustments.
% \end{itemize}

% \end{frame}


\subsection{Constraint Qualification}

\begin{frame}{Why Constraint Qualification Matters}

\textbf{Goal:} ensure that Lagrange multiplier conditions are \emph{valid} and \emph{sufficiently descriptive} of optimal solutions.
\pause
\begin{block}{Problem}
\[
\min f(x) \quad 
\text{s.t.}\quad g_i(x)=0,\; i=1,\ldots,m.
\]
\end{block}
\pause
\begin{itemize}
  \item Lagrange multipliers rely on the fact that the constraint surface is smooth.
  \item At an optimum, the gradients \(\nabla g_i(x^\star)\) must meaningfully describe the geometry of feasible directions.
  \item If they are \textbf{degenerate} (e.g., zero, colinear), the KKT/FOC conditions may fail or produce spurious solutions.
\end{itemize}

\textcolor{BootcampBlueDark}{In short: constraint qualification guarantees that the constraint surface is “regular’’ enough for KKT/FOC to hold.}

\end{frame}

\begin{frame}{Linear Independence Constraint Qualification (LICQ)}

\begin{definitionblock}{Regular (or nondegenerate) point}
A feasible point \(x^\star\) is \textbf{regular} if the gradients
\vspace{-3mm}
\[
\nabla g_1(x^\star),\;\nabla g_2(x^\star),\ldots,\nabla g_m(x^\star)
\vspace{-3mm}
\]
are \textbf{linearly independent}.
\end{definitionblock}
\pause
\textbf{LICQ:}  
If \(x^\star\) is regular, we say that the \emph{linear independence constraint qualification} holds.

\begin{itemize}
  \item Guarantees the existence of Lagrange multipliers.
  \item Ensures \(\nabla f(x^\star)\) can be expressed uniquely as  
        \[
        \nabla f(x^\star)=\sum_{i=1}^m \lambda_i^\star\,\nabla g_i(x^\star).
        \]
  \item Avoids flat regions or “corners’’ in the constraint manifold.
\end{itemize}

\end{frame}

\begin{frame}{Geometric Interpretation}

\begin{itemize}
  \item Each constraint \(g_i(x)=0\) defines a smooth surface.
  \pause
  \item The gradient \(\nabla g_i(x)\) is perpendicular to this surface.
  \pause
  \item LICQ requires these normals to span a full \(m\)-dimensional subspace.
\end{itemize}
  \pause
\begin{center}
\begin{tikzpicture}[scale=1]

  % Axes
  \draw[->] (-0.2,0) -- (3.5,0) node[right] {$x$};
  \draw[->] (0,-0.2) -- (0,3.2) node[above] {$y$};

  % Point optimum
  \coordinate (xstar) at (1.7,1.6);

  % Contrainte g1(x)=0 (bleu)
  \draw[thick, BootcampBlueDark]
    (0.5,2.4)
      .. controls (1.0,1.3) and (1.3,1.6) ..
    (xstar)
      .. controls (2.2,2.0) and (2.7,2.6) ..
    (3.1,2.8);
  \node[BootcampBlueDark] at (2.9,2.4) {$g_1(x)=0$};

  % Contrainte g2(x)=0 (rouge)
  \draw[thick, Crimson]
    (0.6,0.6)
      .. controls (1.1,1.0) and (1.3,1.6) ..
    (xstar)
      .. controls (2.0,2.2) and (2.6,2.4) ..
    (3.0,2.0);
  \node[Crimson] at (2.9,1.0) {$g_2(x)=0$};

  % Point d'intersection
  \fill (xstar) circle (2pt);
  \node[above right] at (xstar) {$x^\star$};

  % Gradient de g1 : normal à g1
  \draw[->, thick, AccentGold]
    (xstar) -- ++(0.9,-0.2)
    node[right] {$\nabla g_1(x^\star)$};

  % Gradient de g2 : normal à g2
  \draw[->, thick, AccentGold]
    (xstar) -- ++(-0.5,0.9)
    node[left] {$\nabla g_2(x^\star)$};

\end{tikzpicture}
\end{center}

\textbf{If the two gradients were colinear, the surfaces would be tangent → LICQ fails.}

\end{frame}

\begin{frame}{When Constraint Qualification Fails}

\textbf{FOC/KKT can break down:}
\begin{itemize}
  \item No Lagrange multipliers may exist.
  \item KKT multipliers exist but do not correspond to a true optimum.
  \item Stationary points of the Lagrangian \emph{miss} the real minimizer.
\end{itemize}
  \pause
\textbf{Example (failure):}
\[
\min f(x)=x 
\quad\text{s.t.}\quad g(x)=x^2=0.
\]
  \pause
\begin{itemize}
  \item Feasible set = \(\{0\}\), but \(\nabla g(0)=0\) so constraint gradient is zero.
  \item LICQ fails.
  \item Lagrange condition: $
    1 + \lambda\cdot 0 = 0 \quad\Rightarrow\quad \text{no solution}.
    $
  \item Yet the true minimizer is clearly \(x^\star=0\).
\end{itemize}

\textcolor{Crimson}{Moral: without CQ, Lagrange FOC \underline{do not} characterize optima.}

\end{frame}

\begin{frame}{Multiple Equality Constraints}

Consider
\[
g(x)=c \quad\Leftrightarrow\quad
\begin{pmatrix}
g_1(x) \\
\vdots \\
g_m(x)
\end{pmatrix}
=
c.
\]
  \pause
The Lagrangian becomes:
\[
\mathcal{L}(x,\lambda)
=f(x)+\lambda^\top(g(x)-c).
\]
  \pause
\textbf{LICQ requirement:}
\[
\nabla g_1(x^\star),\ldots,\nabla g_m(x^\star)
\quad \text{linearly independent.}
\]

\begin{itemize}
  \item Hessian of the Lagrangian must be restricted to feasible directions.
  \item Second-order conditions rely critically on LICQ.
  \item Guarantees that the constraint manifold is a smooth \((n-m)\)-dimensional surface.
\end{itemize}

\end{frame}

\subsection{Inequality Constraints (KKT Conditions)}

\begin{frame}{Optimization with Inequality Constraints}

\textbf{General problem:}
\[
\min_{x\in\mathbb{R}^n} f(x)
\quad\text{s.t.}\quad 
g_i(x)\le 0,\quad i=1,\ldots,m.
\]
  \pause
\begin{itemize}
  \item \(g_i(x)\le 0\) defines a feasible region (possibly with boundaries).
  \item At the optimum, some constraints may bind (\(g_i(x^\star)=0\)).
  \item Others may be slack (\(g_i(x^\star)<0\)).
\end{itemize}
  \pause
\textbf{Idea:}  
Interior points behave like unconstrained minima;  
\emph{boundary points require multipliers and complementarity relations.}

\end{frame}

\begin{frame}{KKT Conditions (Necessary Under CQ)}

\textbf{Lagrangian:}
\[
\mathcal{L}(x,\lambda)
=
f(x)+\sum_{i=1}^m \lambda_i g_i(x),
\qquad \lambda_i\ge 0.
\]
  \pause
If constraint qualification holds (e.g. LICQ), any local minimum \(x^\star\) satisfies:
  \pause
\begin{block}{KKT Conditions}
\[
\textbf{(1) Stationarity: } 
\nabla f(x^\star)+\sum_{i=1}^m\lambda_i^\star\nabla g_i(x^\star)=0
\]
\[
\textbf{(2) Primal feasibility: }
g_i(x^\star)\le 0
\]
\[
\textbf{(3) Dual feasibility: }
\lambda_i^\star\ge 0
\]
\[
\textbf{(4) Complementary slackness: } 
\lambda_i^\star g_i(x^\star)=0.
\]
\end{block}
  \pause
\end{frame}

\begin{frame}{Complementary Slackness: Interpretation}

\[
\lambda_i^\star g_i(x^\star)=0 \quad \forall i.
\]
  \pause
\begin{itemize}
  \item If constraint is \textbf{slack}  $
        g_i(x^\star)<0,
        $
        it does not affect optimality so \(\lambda_i^\star=0\).
          \pause
  \item If constraint is \textbf{active / binding}  
        $
        g_i(x^\star)=0,
        $
        the multiplier may be strictly positive.
          \pause
  \item \(\lambda_i^\star\) measures the \textbf{shadow price} of relaxing the constraint.
\end{itemize}

\begin{center}
\begin{tikzpicture}[scale=1]
  \draw[->] (-0.2,0) -- (4.2,0) node[right] {$x$};
  \draw[->] (0,-0.2) -- (0,3.2) node[above] {$f(x)$};

  % feasible region shading
  \draw[fill=gray!20] (2,0) -- (4,0) -- (4,3) -- (2,3) -- cycle;

  \node at (3,2.8) {feasible};
  \draw[dashed] (2,0) -- (2,3) node[midway,left] {$g(x)=0$};

  % function
  \draw[thick, BootcampBlueDark] plot[smooth] coordinates {(0.5,1.5) (1,1) (1.5,0.8) (2,0.7) (2.5,0.8) (3,1.2) (3.5,2)};

  % optimum at boundary
  \fill (2,0.7) circle (2pt) node[below right] {$x^\star$ (binding)};
\end{tikzpicture}
\end{center}

\end{frame}

\begin{frame}{Example (1D): Boundary Optimum}

\[
\min_{x} \ f(x)=x^2-4x
\quad \text{s.t.} \quad g(x)=1-x \le 0.
\]
\pause
\textbf{Unconstrained minimum:}  
\(f'(x)=2x-4=0 \Rightarrow x=2\).  
But feasible set is \(x \ge 1\).
\pause
\textbf{Check feasibility:}  
Constraint active at optimum?  
Yes: best feasible point is \(x^\star=2\) (interior).
\pause
\textbf{KKT:}
\[
\nabla f(2)=0,\quad g(2)= -1<0\Rightarrow \lambda^\star=0.
\]

Thus:
\[
x^\star=2,\qquad \lambda^\star=0.
\]

\textbf{Moral:}  
Interior optimum imply that multiplier vanishes.

\end{frame}

\begin{frame}{Example (1D): Binding Constraint}

\[
\min_x f(x)=x^2
\quad\text{s.t.}\quad x\ge 1 \ \Leftrightarrow\ g(x)=1-x\le 0.
\]

Unconstrained minimum: \(x=0\), but infeasible.  
Feasible minimum at boundary \(x^\star=1\).
\pause
\textbf{KKT:}

Stationarity:
\[
\nabla f(1) + \lambda^\star(-1) = 0
\quad\Rightarrow\quad
2 - \lambda^\star = 0
\; \Rightarrow\;
\lambda^\star = 2>0.
\]

Complementary slackness:
\(
\lambda^\star g(1)=2\cdot 0=0.
\)
\pause
\textbf{Conclusion:}  
\[
x^\star=1,\quad 
\lambda^\star=2.
\]

The constraint is active → multiplier positive.

\end{frame}

\begin{frame}{Example (2D): Boundary Minimum}

\[
\min_{x,y} f(x,y)=(x-1)^2 + (y-2)^2
\quad\text{s.t.}\quad g(x,y)=x+y-1\le 0.
\]

Unconstrained optimum: \((1,2)\) but infeasible  
(\(1+2-1=2>0\)).
\pause
\textbf{Closest feasible point lies on the line}  $
x+y=1.
$
$$
\begin{pmatrix}
2(x-1)\\[2pt]
2(y-2)
\end{pmatrix}
+ \lambda 
\begin{pmatrix}
1\\ 1
\end{pmatrix}
=0.
$$
\pause
Constraint:
$
x+y=1.
$
Solve:
$
x^\star=0,\quad y^\star=1,\quad \lambda^\star= -2.
$
\pause
Dual feasibility requires \(\lambda^\star\ge 0\):  
But note: constraint is written \(g=x+y-1\le 0\).  
Gradient is \(\nabla g=(1,1)\).  
Sign consistency gives:
\[
\lambda^\star = 2 >0.
\]

\textbf{Boundary optimum with positive multiplier.}

\end{frame}

\begin{frame}{When KKT Conditions Fail}

If the constraint gradients are linearly dependent or zero,  
\textbf{KKT may not hold} even at true minima.
\pause
\textbf{Example:}
\[
\min f(x)=x 
\quad\text{s.t.}\quad g(x)=x^2\le 0.
\]

Feasible set: \(\{0\}\).  
True minimum is \(x^\star=0\).  

But:
\[
\nabla g(0)=0 \quad\Rightarrow\quad \text{LICQ fails.}
\]

Trying KKT:
\[
1 + \lambda \cdot 0 = 0 \Rightarrow \text{no solution}.
\]

\textbf{Conclusion:}  
KKT does NOT detect the true optimum →  
\textcolor{Crimson}{constraint qualification is essential.}

\end{frame}



% % ======================================
% % SECTION 4 — Applications
% % ======================================
% \section{Applications}

% \begin{frame}{Outline}
%   \tableofcontents[currentsection, hideothersubsections]
% \end{frame}

% \subsection{Mean-Variance Portfolio Optimization}
% \begin{frame}{Markowitz Portfolio Problem}
% \[
% \min_{w} \;\; w^\top \Sigma\, w
% \qquad\text{s.t.}\qquad 
% w^\top \mathbf{1}=1,\,\;
% \mathbb{E}[R]^\top w = \mu.
% \]
% \begin{itemize}
%   \item Quadratic objective + linear constraints.
%   \item Closed-form solution via Lagrange multipliers.
% \end{itemize}
% % TODO: sketch mean-variance frontier
% \end{frame}

% \subsection{Budget Constraints \& Shadow Prices}
% \begin{frame}{Shadow Prices}
% \begin{itemize}
%   \item Lagrange multiplier represents marginal value of relaxing a constraint.
%   \item Dual interpretation: price of the constrained resource.
% \end{itemize}
% % TODO: add simple consumption/budget example
% \end{frame}

% \subsection{Calibration \& Least Squares}
% \begin{frame}{Least-Squares Estimation}
% \[
% \min_x \|Ax-b\|_2^2
% \]
% \begin{itemize}
%   \item Leads to normal equations \(A^\top A x = A^\top b\).
%   \item Core tool for calibration and curve fitting.
% \end{itemize}
% \end{frame}

% \subsection{Linear Regression}
% \begin{frame}{Linear Regression as Optimization}
% \begin{itemize}
%   \item Model: \(y = X\beta + \varepsilon\).
%   \item OLS estimator: \(\hat\beta = (X^\top X)^{-1} X^\top y\).
%   \item Minimizes sum of squared residuals = least-squares optimization.
% \end{itemize}
% \end{frame}

%------------------------------------------------
\section{Nonsmooth convex optimization}
%------------------------------------------------

\begin{frame}{Why nonsmooth convex optimization?}
\begin{itemize}
    \item Many important convex objectives are not differentiable everywhere:
    \begin{itemize}
        \item $\ell_1$ norm $\|x\|_1$, absolute value $|x|$
        \item hinge loss, ReLU, max of affine functions
    \end{itemize}
    \pause
    \item Classical gradient-based methods use $\nabla f(x)$
    \begin{itemize}
        \item $\Rightarrow$ not directly applicable at nondifferentiable points
    \end{itemize}
    \pause
    \item Goal: extend first-order theory (gradient, optimality conditions, algorithms)
    \begin{itemize}
        \item to convex functions that are not smooth
    \end{itemize}
\end{itemize}
\end{frame}

%------------------ SUBGRADIENT ------------------

\begin{frame}{Subgradients and subdifferential}
\begin{block}{Definition (subgradient)}
Let $f:\mathbb{R}^n \to \mathbb{R}$ be convex.  
A vector $g \in \mathbb{R}^n$ is a \textbf{subgradient} of $f$ at $x$ if
\[
f(y) \ge f(x) + g^\top (y-x), \quad \forall y \in \mathbb{R}^n.
\]
\end{block}
\pause

\begin{block}{Subdifferential}
The \textbf{subdifferential} of $f$ at $x$ is the set
\[
\partial f(x) = \{ g \in \mathbb{R}^n : g \text{ is a subgradient of } f \text{ at } x \}.
\]
\end{block}
\pause
\begin{itemize}
    \item If $f$ is differentiable at $x$, then $\partial f(x) = \{ \nabla f(x)\}$.
    \item At a “kink”, $\partial f(x)$ is typically a nonempty convex set.
\end{itemize}
\end{frame}

\begin{frame}{Subgradients: geometric intuition}
\begin{center}
\begin{tikzpicture}[scale=1.1]

  % Axes
  \draw[->] (-0.5,0) -- (4.5,0) node[right] {$x$};
  \draw[->] (0,-0.5) -- (0,3.0) node[above] {$f(x)$};

  % Convex piecewise-linear function with a kink at x*
  \draw[thick]
    (0.2,2.3) -- (2.0,1.0) -- (4.0,2.4);

  % Point of nondifferentiability (kink)
  \coordinate (xk) at (2.0,1.0);
  \fill (xk) circle (2pt);
  \node[below right] at (xk) {$x^\star$};

  % Supporting lines (subgradients) through x*
  % slopes m = -0.6, 0, 0.6 (all ≤ right slope and ≥ left slope)
  \draw[AccentGold,thick] (0.0,2.2) -- (4.2,-0.32);
  \draw[AccentGold,thick] (0.0,1.0) -- (4.2,1.0);
  \draw[AccentGold,thick] (0.0,-0.2) -- (4.2,2.32);

  \node[AccentGold] at (3.9,2.6) {supporting lines};

\end{tikzpicture}
\end{center}

\begin{itemize}
\item At $x^\star$: many supporting lines $\Rightarrow$ $\partial f(x^\star)$ is a set.
\item At a smooth point $x_d$: unique tangent $\Rightarrow$ $\partial f(x_d)=\{\nabla f(x_d)\}$.
\end{itemize}

\end{frame}

\begin{frame}{Examples of subdifferentials}
\begin{itemize}
    \item Absolute value in 1D:
    \[
    f(t) = |t|
    \]
    \[
    \partial f(t) =
    \begin{cases}
        \{1\}  & \text{if } t>0,\\
        [-1,1] & \text{if } t=0,\\
        \{-1\} & \text{if } t<0.
    \end{cases}
    \]
    \item $\ell_1$ norm in $\mathbb{R}^n$:
    \[
    f(x) = \|x\|_1 = \sum_{i=1}^n |x_i|.
    \]
    A subgradient $g\in\partial f(x)$ satisfies
    \[
    g_i =
    \begin{cases}
        \operatorname{sign}(x_i) & \text{if } x_i\neq 0,\\
        \gamma_i, \ |\gamma_i|\le 1 & \text{if } x_i=0.
    \end{cases}
    \]
\end{itemize}
\end{frame}

%------------------ OPTIMALITY ------------------

\begin{frame}{Optimality condition (nonsmooth case)}
We consider the unconstrained problem
\[
\min_{x\in\mathbb{R}^n} f(x), \qquad f\ \text{convex}.
\]
\pause
\begin{block}{First-order optimality condition}
A point $x^\star$ is a (global) minimizer of $f$ if and only if
\[
0 \in \partial f(x^\star).
\]
\end{block}
\pause
\begin{itemize}
    \item Generalizes $\nabla f(x^\star)=0$ in the smooth case.
    \item Because of convexity, any point satisfying $0\in \partial f(x^\star)$ is globally optimal.
\end{itemize}
\end{frame}

%------------------ SUBGRADIENT METHOD ------------------

\begin{frame}{Subgradient method}
We aim to solve
\[
\min_x f(x), \quad f \text{ convex, but not necessarily differentiable}.
\]
\pause
\begin{block}{Algorithm (subgradient method)}
Choose $x_0$ and a stepsize sequence $(\alpha_k)_{k\ge 0}$.  
For $k=0,1,2,\dots$
\[
g_k \in \partial f(x_k),
\qquad
x_{k+1} = x_k - \alpha_k g_k.
\]
\end{block}
\pause
\begin{itemize}
    \item Very general: only requires access to subgradients.
    \item Function values do \textbf{not} necessarily decrease at each step.
    \item Convergence is typically slower than in the smooth gradient case.
\end{itemize}
\end{frame}

\begin{frame}{Convergence (idea)}
Assume:
\begin{itemize}
    \item $f$ is convex and has a minimizer $x^\star$,
    \item subgradients are bounded: $\|g_k\|\le G$,
    \item stepsizes satisfy
    \[
    \alpha_k >0,\quad \sum_{k=0}^\infty \alpha_k = \infty,\quad
    \sum_{k=0}^\infty \alpha_k^2 < \infty.
    \]
\end{itemize}
\pause
\begin{block}{Result (informal)}
Then the sequence of function values satisfies
\[
f(x_k) \to f(x^\star) \quad \text{as } k\to\infty.
\]
\end{block}
\pause
Typical choice:
\[
\alpha_k = \frac{c}{\sqrt{k+1}}.
\]
\end{frame}

%------------------ PROX OPERATOR ------------------

\begin{frame}{Proximal operator}
Let $g:\mathbb{R}^n\to\mathbb{R}\cup\{+\infty\}$ be proper, convex and lower semicontinuous.
\pause
\begin{block}{Definition (proximal operator)}
For $\lambda>0$ and $v\in\mathbb{R}^n$,
\[
\operatorname{prox}_{\lambda g}(v)
=
\arg\min_{x\in\mathbb{R}^n}
\left\{
g(x) + \frac{1}{2\lambda}\|x - v\|^2
\right\}.
\]
\end{block}
\pause
\begin{itemize}
    \item Combines a nonsmooth term $g$ with a quadratic regularization.
    \item Yields a point that balances being close to $v$ and having small $g(x)$.
\end{itemize}
\pause
\end{frame}

\begin{frame}{Examples of proximal operators}
\begin{itemize}
    \item Indicator of a closed convex set $C$:
    \[
    g(x) = \iota_C(x) =
    \begin{cases}
       0 & \text{if } x\in C,\\
       +\infty & \text{otherwise}.
    \end{cases}
    \]
    \pause
    Then
    \[
    \operatorname{prox}_{\lambda g}(v) = P_C(v)
    \]
    is the Euclidean projection onto $C$.
    \pause
    \item $\ell_1$ norm:
    \[
    g(x)=\|x\|_1 = \sum_i |x_i|.
    \]
    The proximal operator is the \textbf{soft-thresholding} operator:
    \[
    (\operatorname{prox}_{\lambda \|\cdot\|_1}(v))_i
    = \operatorname{sign}(v_i)\max\{|v_i| - \lambda,0\}.
    \]
\end{itemize}
\end{frame}

%------------------ PROX GRADIENT ------------------

\begin{frame}{Use of prox: proximal gradient method}
Consider a composite problem
\[
\min_x F(x) := f(x) + g(x),
\]
where
\begin{itemize}
    \item $f$ is smooth (differentiable with Lipschitz gradient),
    \item $g$ is convex, possibly nonsmooth (e.g. $\ell_1$ penalty, indicator of a set).
\end{itemize}

\begin{block}{Proximal gradient iteration}
Given $x_k$ and a stepsize $\lambda>0$,
\[
x_{k+1}
=
\operatorname{prox}_{\lambda g}\Big(x_k - \lambda \nabla f(x_k)\Big).
\]
\end{block}

\begin{itemize}
    \item Gradient step for the smooth part $f$,
    \item Proximal step for the nonsmooth part $g$.
\end{itemize}
\end{frame}

\begin{frame}{Applications of proximal gradient}
\begin{itemize}
    \item \textbf{Sparse regression} (LASSO):
    \[
    \min_x \frac{1}{2}\|Ax-b\|^2 + \lambda \|x\|_1,
    \]
    where $f(x)=\frac{1}{2}\|Ax-b\|^2$ and $g(x)=\lambda\|x\|_1$.
    \pause
    \item \textbf{Constrained problems}:
    \[
    \min_{x\in C} f(x)
    \quad \Longleftrightarrow \quad
    \min_x f(x) + \iota_C(x),
    \]
    leading to
    \[
    x_{k+1}=P_C\big(x_k - \lambda \nabla f(x_k)\big)
    \]
    (projected gradient method).
    \item Numerous problems in signal processing, imaging, and machine learning.
\end{itemize}
\end{frame}

%------------------------------------------------
\section{Algorithmic methods for optimization}
%------------------------------------------------

\begin{frame}{Optimization problems: classification}
We consider problems of the form
\[
\min_{x} \; F(x)
\]
possibly subject to constraints.
\pause
\begin{itemize}
    \item \textbf{unconstrained} vs \textbf{constrained}
    \item \textbf{smooth} vs \textbf{nonsmooth}
\end{itemize}

Different algorithmic strategies are appropriate depending on the case.
\end{frame}

%================= UNCONSTRAINED SMOOTH =================

\begin{frame}{Unconstrained smooth optimization}
We first assume:
\begin{itemize}
    \item $F$ differentiable
    \item $\nabla F$ is Lipschitz continuous
\end{itemize}
\pause
\begin{block}{Gradient descent}
\[
x_{k+1} = x_k - \alpha_k \nabla F(x_k)
\]
\end{block}
\pause
\begin{block}{Newton’s method}
\[
x_{k+1} = x_k - \nabla^2 F(x_k)^{-1}\nabla F(x_k)
\]
\end{block}

\begin{itemize}
    \item gradient descent: cheap per iteration, slower convergence
    \item Newton: expensive per iteration, fast local convergence
    \item quasi–Newton (BFGS, L-BFGS): approximation of Hessian
\end{itemize}
\end{frame}

%================= UNCONSTRAINED NONSMOOTH =================

\begin{frame}{Unconstrained nonsmooth optimization}
Now $F$ is convex but may not be differentiable.
\pause
\begin{itemize}
    \item absolute value
    \item $\ell_1$ norm
    \item piecewise-linear convex functions
\end{itemize}
\pause
\begin{block}{Subgradient method}
\[
x_{k+1}=x_k-\alpha_k g_k,
\qquad g_k\in \partial F(x_k)
\]
\end{block}
\vspace{-5mm}
\begin{itemize}
    \item most general method
    \item convergence in function value (slow rates)
\end{itemize}
\pause
\begin{block}{Proximal methods (composite problems)}
\[
\min_x \; f(x)+g(x)
\]
with $f$ smooth and $g$ convex: $
x_{k+1} = \operatorname{prox}_{\lambda g}\left(x_k-\lambda\nabla f(x_k)\right)
$
\end{block}
\end{frame}

%================= CONSTRAINED SMOOTH =================

\begin{frame}{Constrained smooth optimization}
We now solve:
\[
\min_{x\in C} f(x)
\]
with $f$ smooth and $C$ convex.
\pause
\begin{block}{Projected gradient method}
\[
x_{k+1}
=
P_C\left(x_k-\alpha\nabla f(x_k)\right)
\]
\end{block}
\pause
\begin{itemize}
    \item gradient step
    \item projection onto feasible set
\end{itemize}

Other classical approaches:
\begin{itemize}
    \item penalty and barrier methods
    \item Lagrange multipliers and KKT conditions
    \item interior–point methods
\end{itemize}
\end{frame}

%------------------------------------------------

\begin{frame}{Lagrangian formulation}
For some convex $f$, consider constrained problems
\[
\min_x f(x) \quad \text{s.t. } \quad Ax = b.
\vspace{-8mm}
\]
\pause
\begin{block}{Lagrangian}
\[
\mathcal{L}(x,\lambda)
=
f(x) + \lambda^\top (Ax-b),
\]
where $\lambda$ are Lagrange multipliers.
\end{block}
\pause
\begin{itemize}
    \item primal variable: $x$
    \item dual variable: $\lambda$
\end{itemize}

Optimization can be done:
\begin{itemize}
    \item on the primal
    \item on the dual
    \item or on both (primal–dual)
\end{itemize}
\end{frame}

%================= DUAL ASCENT =================

\begin{frame}{Dual problem and dual ascent}
The dual function is
\[
d(\lambda)=\inf_x \mathcal{L}(x,\lambda),
\]
and the dual optimization problem is
\[
\max_{\lambda} \; d(\lambda).
\]
\pause
\begin{block}{Dual ascent method}
Iterate:
\[
x_{k+1} = \arg\min_x \mathcal{L}(x,\lambda_k),
\]
\[
\lambda_{k+1}
=
\lambda_k + \alpha_k (Ax_{k+1}-b).
\]
\end{block}

\begin{itemize}
    \item alternating minimization in $x$ and ascent in $\lambda$
    \item works best when $\mathcal{L}$ is easily minimized in $x$
\end{itemize}
\end{frame}

%================= DUAL SUBGRADIENT =================

\begin{frame}{Subgradient method on the dual}
Even when $d(\lambda)$ is not differentiable, it is concave.

\begin{block}{Dual subgradient iteration}
\[
\lambda_{k+1}
=
\lambda_k
+
\alpha_k (Ax_k - b),
\]
where
\[
x_k \in \arg\min_x \mathcal{L}(x,\lambda_k).
\]
\end{block}
\pause
\begin{itemize}
    \item $Ax_k-b$ is a subgradient of $d(\lambda)$
    \item convergence under standard diminishing step size conditions
\end{itemize}
\end{frame}

%================= METHOD OF MULTIPLIERS =================

\begin{frame}{Augmented Lagrangian / method of multipliers}
To improve convergence, add a quadratic penalty term:

\[
\mathcal{L}_\rho(x,\lambda)
=
f(x)+\lambda^\top (Ax-b)
+
\frac{\rho}{2}\|Ax-b\|^2.
\]
\pause
\begin{block}{Method of multipliers}
\[
x_{k+1}
=
\arg\min_x \mathcal{L}_\rho(x,\lambda_k),
\]
\[
\lambda_{k+1}
=
\lambda_k + \rho(Ax_{k+1}-b).
\]
\end{block}
\pause
\begin{itemize}
    \item less sensitive to ill-conditioning than pure dual ascent
    \item strong convergence properties for convex problems
\end{itemize}
\end{frame}

%================= PRIMAL-DUAL =================

\begin{frame}{Primal–dual methods}
Consider simultaneously
\[
\min_x f(x)\quad\text{and}\quad\max_\lambda d(\lambda).
\]
\pause
\begin{itemize}
    \item update $x$ (primal step)
    \item update $\lambda$ (dual step)
\end{itemize}
\pause
Examples include:
\begin{itemize}
    \item Arrow–Hurwicz method
    \item Chambolle–Pock algorithm
\end{itemize}
\pause
Primal–dual methods are particularly useful for
\begin{itemize}
    \item saddle-point problems
    \item variational imaging
    \item constrained convex optimization
\end{itemize}
\end{frame}

%================= CONSTRAINED NONSMOOTH =================

\begin{frame}{Constrained nonsmooth optimization}
General form:
\[
\min_{x\in C} f(x) + g(x)
\]
where $g$ may encode nonsmooth regularization and/or constraints.
\pause
\begin{itemize}
    \item indicator functions of sets
    \item sparsity-promoting penalties
\end{itemize}

\begin{block}{Projected subgradient}
\[
x_{k+1}
=
P_C\left(x_k-\alpha_k g_k\right),
\qquad g_k\in \partial f(x_k)
\]
\end{block}
\pause
\begin{block}{Augmented Lagrangian and ADMM}
Handle constraints and nonsmooth terms by splitting:
\[
\min_{x,z} f(x)+g(z)
\quad \text{s.t. } Ax+Bz=c
\]
Alternating updates on $x$ and $z$.
\end{block}
\end{frame}

%================= SUMMARY TABLE =================

\begin{frame}{Summary by case (I)}

\begin{itemize}

    \item \textbf{Unconstrained + smooth}  
    \begin{itemize}
        \item gradient descent
        \item Newton and quasi–Newton methods
        \item accelerated first–order methods
    \end{itemize}

    \item \textbf{Unconstrained + nonsmooth}  
    \begin{itemize}
        \item subgradient method
        \item proximal gradient methods
        \item bundle / cutting–plane methods
    \end{itemize}

    \item \textbf{Constrained + smooth}  
    \begin{itemize}
        \item projected gradient methods
        \item penalty and barrier methods
        \item interior–point methods
    \end{itemize}

\end{itemize}

\end{frame}

\begin{frame}{Summary by case (II)}

\begin{itemize}

    \item \textbf{Constrained + nonsmooth}  
    \begin{itemize}
        \item projected subgradient methods
        \item proximal–projection methods
        \item augmented Lagrangian methods
        \item ADMM (splitting and decomposition)
    \end{itemize}

    \item \textbf{Lagrangian / dual methods}  
    \begin{itemize}
        \item dual ascent and dual subgradient methods
        \item method of multipliers
        \item augmented Lagrangian
        \item primal–dual algorithms
    \end{itemize}

\end{itemize}

\bigskip

Key messages:
\[
\text{proximal operator = nonsmooth analogue of projection}
\]
\[
\text{augmented Lagrangian = Lagrangian + quadratic penalty}
\]

\end{frame}


\end{document}